\chapter{Introduction} \label{intro}

Virtual terrains appear in a wide variety of applications: from entertainment
to physical simulation or science visualization, a lot of virtual scenes 
are set on a digital landscape. Despite the ubiquity of virtual terrains 
on 3D scenes, the complex and fractal structure of natural landscapes still
brings multiple challenges for the field of synthetic terrain generation.

\vspace{10pt}

Most virtual terrains are represented using 2D heightfields which, despite 
being highly convenient when generating, storing and manipulating terrain
models, have a major drawback: heightfields can only represent 2.5D 
models. Because each point in a 2D plane has a set elevation and every 
point between the ground and the given elevation is considered as solid 
(and the terrain interior), we can not represent features like caves, 
arches or overhangs, where not all points between the ground and the 
maximum elevation are solid. Additionally, these representations of terrains 
also fail to reproduce complex geometry commonly found in rocky terrains, 
like sharp ridges, loose blocks or steep cliffs.

\vspace{10pt}

Many of these complex structures that can not be easily captured by 2D 
heightfields, arise from different types of weathering events, like percolation
and freeze-thaw cycles. To address this limitations on virtual terrain models,
our goal in this thesis is to study the mechanical weathering of rocky 
terrains to enhance models given from 2D heightfields and add the complex features 
that these fields fail to capture. 

\vspace{10pt}

Terrain enhancement through weathering simulation not only improves the 
realism and quality of the models, adding the aforementioned complex 
features, but it can also be applied dynamically to simulate the changes 
a terrain goes through over multiple years.

\vspace{10pt}

This project is built on a previous project on this same topic \cite{10.2312:ceig.20241139}. 
Using the same algorithm as a base, we worked on addressing some of its 
limitations and improving both its quality and efficiency.
\section{Contributions}

As we just mentioned, this project is heavily based on a previous one, and
our proposed algorithm is built on the method they proposed. Despite that, 
in this thesis we will provide 
multiple improvements to the base algorithm to address some of its limitations. 

\vspace{10pt}

The main point we wanted to address of the previous project was the lack of
a physically based approach to the algorithm, as the proposed algorithm
made a lot of simplifications and ignored multiple physical interactions 
which made the method less accurate and realistic. However, during this 
thesis we also found other areas that could be improved and expanded, 
which also took a significant ammount of effort to solve.

\vspace{10pt}

The limitations of the previous project and our improvements will be 
discussed at length during this thesis, but our main contributions can
be summarized in two large blocks: improvements to the decomposition of 
the terrain into cells and
improvements on the quality of the algorithm by including a mechanical 
model to the simulation.  

\section{Set-up and tools}

Another difference between the project we base this thesis on, albeit less
significant, resides in the tools used. In the project that presented the algorithm we 
are extending, a blender plugin was implemented to show the method, however, 
here we have implemented our erosion simulation algorithm using C++ and we showed the 
results in a QT application using OpenGL for rendering.

\vspace{10pt}

To make the application, we reused the code provided by an article on 
terrain descriptors \cite{Argudo2025descriptors}, as it already 
implemented many basic functionalities to work with virtual terrains. We
then modified and expanded the application implementing our algorithm for 
mechanical weathering as well as implementing other functionalities to aid in the 
visualization and debugging of our method. Some of the changes that we 
implemented modify the basic functionalities of the program, so 
our project should not be seen as an extension to said application.

\vspace{10pt}

Finally, we used the C++ library Voro++ \cite{osti_code-54674} to generate 
voronoi decompositions, a core aspect of our erosion algorithm.